
In this section, we will review the ABC framework and highlight its limitations within the hierarchical model framework. Subsequently, we will introduce our permutation-based method, permABC, which aims to address these limitations.

\subsection{Introduction to ABC}


The standard rejection ABC algorithm, often referred to as \emph{Vanilla ABC}, proceeds by drawing parameters from the prior distribution, $\param \sim \pi(\cdot)$, and then simulating synthetic data from the likelihood, $\simdata \sim f(\cdot \mid \param)$. The pair $(\param, \simdata)$ is accepted if the summary statistics of the simulated data are sufficiently close to those of the observed data, that is, if $d(\eta(\obsdata), \eta(\simdata)) \leq \varepsilon$ for some tolerance level $\varepsilon > 0$ and some summary statistic $\eta$. This procedure yields samples from the approximate joint distribution $\pi_{\varepsilon}\{\param,\simdata \mid \eta(\obsdata)\}$; we marginalize out $\simdata$ to get the ABC posterior

\begin{equation}
    \label{eqabc}
    \begin{split}
        \pi_{\varepsilon}(\param \mid \eta(\obsdata)) &\propto
        \int_{\dataspace^\Clust} \pi_{\varepsilon}(\param, \simdata \mid \eta(\obsdata)) \, d\simdata \\
        &\propto \pi(\param) \int_{A_{\obsdata,\varepsilon}} f(\simdata \mid \param) \, d\simdata,
    \end{split}
\end{equation}
where $A_{\obsdata,\varepsilon} = \{\simdata \in \dataspace^\Clust \mid d(\eta(\obsdata), \eta(\simdata)) \leq \varepsilon\}$ denotes the acceptance region.

Combining an informative summary statistic $\eta$ with a small tolerance $\varepsilon$ typically leads to a good approximation of the true posterior distribution, so that $\pi_{\varepsilon}(\param \mid \eta(\obsdata)) \approx \pi(\param \mid \obsdata)$. This approximate posterior—often called a \emph{pseudo-posterior}—is obtained by replacing the true likelihood $f(\obsdata \mid \param)$ with the integral of the likelihood over the neighborhood $A_{\obsdata,\varepsilon}$ of the observed data, effectively using it as a surrogate for the intractable likelihood.

Throughout this work, we assume for clarity that the data are compared directly, i.e., $\eta = \mathrm{Id}$. Under this assumption, the acceptance region simplifies to the $\varepsilon$-ball $B_{\varepsilon}(\obsdata)$ centered at $\obsdata$ with respect to the distance $d$. Nonetheless, the proposed method naturally extends to the use of low-dimensional summary statistics.


In our numerical examples, we take $\dataspace = \mathbb R^n$ and use the distance $d(\cdot,\cdot)$ between two datasets in \( \dataspace^K \), each consisting of \( K \times n \) observations, defined as follows: 
\begin{equation}
    \label{eqdistance}
    d(\obsdata, \simdata) = \sqrt{\sum_{k=1}^K w_k^2 \| \obsdata_k - \simdata_k \|_2^2},
\end{equation}
where \( \| \cdot \|_2 \) denotes the Euclidean norm, and the weights \( w_k \) are used to control the relative influence of each compartment in the distance computation. The weight vector \( w_{1:K} \) can either be manually specified by the user or selected adaptively to prevent any single compartment from dominating the distance, as suggested by \citet{prangle2017adapting}.


This naive ABC algorithm faces several limitations, including poor scalability in the parameter space and an exponentially increasing computational cost due to the declining acceptance rate as \( \varepsilon \) decreases. In the context of the hierarchical model of Figure \ref{fig:hier}, the acceptance rate further diminishes as the number of compartments \( K \) increases. Simulating a dataset \( \simdata \) from the prior predictive distribution that matches all compartments of \( \obsdata \) within a small tolerance \( \varepsilon \) can be particularly challenging. 


\subsection{Motivations of permABC}

A common approach in ABC to increase the acceptance rate is to exploit the exchangeability of the compartments through a permutation-invariant summary statistic, such as the sorted vector. For instance, when there is a one-dimensional observation (or statistic) per compartment, $\obsdata \in \mathbb{R}^K$, it is natural to use the ordering $\eta(\obsdata)=\mathrm{sort}(\obsdata)$ as a summary. This replaces the $K!$ possible labelings by a canonical representation, thereby simplifying the comparison between datasets.

Moreover, in the univariate case, sorting is known to minimize the distance over all permutations of the simulated data. Specifically, \(
d\bigl(\mathrm{sort}(\obsdata),\mathrm{sort}(\simdata)\bigr)
=
\min_{\sigma\in\mathcal{S}_K} d(\obsdata,\simdata_\sigma),
\)
where $\mathcal{S}_K$ denotes the set of all permutations of $\{1,\dots,K\}$. Equivalently, in this one-dimensional setting, comparing sorted vectors amounts to computing the Wasserstein distance between the empirical measures of the compartments, which connects our construction to Wasserstein ABC (WABC) \citep{bernton2019approximate}.

However, this equivalence should not obscure an important difference in inferential target. WABC is designed to compare samples viewed as unordered point clouds, and is therefore primarily suited to settings where inference concerns global features of the data distribution. By contrast, our setting is hierarchical: inference concerns not only the global parameter \( \globparam \), but also the local parameters \( \locparam_{1:K} \), at least up to relabelling. In particular, WABC is not designed for hierarchical models with compartment-specific local parameters, since it does not keep track of which simulated compartment is associated with which local parameter.

In higher dimension, the scalar sorting argument no longer applies. When each $\obsdata_k$ is a vector in $\mathbb{R}^n$ with $n>1$, there is no canonical ordering that is both deterministic and permutation-optimal. We therefore solve the matching problem directly, which leads to the acceptance criterion
\begin{equation}
    \label{eqopti}
    \min_{\sigma \in \mathcal{S}_K} d(\obsdata,\simdata_\sigma)\le \varepsilon.
\end{equation}

This condition means that the simulated dataset is accepted whenever it matches the observed dataset up to a permutation of the compartments. It can therefore yield a substantial increase in acceptance probability, heuristically of order $K!$ when $\varepsilon$ is small and the neighborhoods associated with distinct permutations are essentially disjoint. The optimization problem in \eqref{eqopti} can be formulated as a linear sum assignment problem, or equivalently a bipartite matching problem, and solved efficiently using a modified Jonker--Volgenant algorithm \citep{jonker1987shortest}, with computational complexity $\mathcal{O}(K^3)$ \citep{crouseImplementing2DRectangular2016a}. Although this remains substantial, it reduces the combinatorial cost from $K!$ to polynomial time and makes problems with up to $K=100$ compartments tractable in practice; see the real-data application in Section~\ref{sec:sir}.

This matching criterion is closely related to optimal transport between the empirical distributions of the compartments. In the multivariate case, \citet{bernton2019approximate} discuss several cheaper approximations to exact transport, including Hilbert-curve sorting, swap-based refinements, and entropically regularized Sinkhorn divergences. These approximations are particularly attractive when the number of points is very large. In our applications, however, $K$ remains only moderately large, while the dimension of each compartment-level observation may be non-negligible. In this regime, solving the assignment problem exactly remains practical and robust. We nevertheless return to this point in Section~\ref{sec:sequential}, where we show that within an SMC scheme a cheaper swap-based refinement, warm-started from the previous iteration's optimal permutation, is sufficient in practice.

However, using \eqref{eqopti} does not yield samples from the pseudo-posterior defined in \eqref{eqabc}, but rather from the following permutation-invariant version:
\begin{equation}
    \label{eqpermabc_tilde}
    \begin{split}
        \Tilde \pi_{\varepsilon} (\param \mid \obsdata)
        &\propto \int_{\dataspace^K} \Tilde \pi_{\varepsilon} (\param,\simdata \mid \obsdata)\, d\simdata \\
        &\propto \pi(\param)\int_{\Tilde A_{\obsdata,\varepsilon}} f(\simdata \mid \param)\, d\simdata,
    \end{split}
\end{equation}
where the acceptance region is
\[
\tilde{A}_{\obsdata,\varepsilon}
=
\left\{
\simdata\in\dataspace^K \;\middle|\;
\min_{\sigma\in\mathcal{S}_K} d(\obsdata,\simdata_\sigma)\le \varepsilon
\right\}
=
\bigcup_{\sigma\in\mathcal{S}_K} B_\varepsilon(\obsdata_\sigma),
\]
with \(B_\varepsilon(\obsdata_\sigma)\) denoting the $\varepsilon$-ball centered at \(\obsdata_\sigma\) for the distance \(d\).

In this formulation, the likelihood \(f(\obsdata\mid\param)\) is approximated by its integral over a neighborhood of all permutations of the observed dataset:
\[
\int_{\tilde A_{\obsdata,\varepsilon}} f(\simdata\mid\param)\, d\simdata
\approx
\sum_{\sigma\in\mathcal{S}_K} f(\obsdata_\sigma\mid\param).
\]

Inference under \eqref{eqpermabc_tilde} is therefore naturally defined only up to a permutation of the local components \( \locparam_{1:K} \), as in WABC \citep{bernton2019approximate}, effectively yielding posterior samples of the form \(
\Tilde{\param}
=
(\globparam,\{\locparam_1,\dots,\locparam_K\}),
\)
where the local parameters are treated as an unordered set. This target may be appropriate when the local parameters are nuisance quantities, but in many hierarchical applications one wishes to recover inference on the ordered local components as well.

To our knowledge, this specific permutation-based formulation of ABC for hierarchical models with local parameters has not been formally studied in the literature. We therefore introduce \emph{permABC}, which exploits permutation matching to improve acceptance while making it possible to recover inference on the full parameter vector \( \param \), including the ordered local parameters.

\subsection{The permABC algorithm}
\label{secpermabc_star}


Restoring identifiability among compartments is essential for enabling inference on the full parameter vector. One natural way to achieve this is by applying a suitable permutation to the accepted parameters, such that the simulated data are optimally aligned with the observed dataset. However, explicitly testing all \( K! \) permutations is computationally intractable.

Fortunately, whenever the acceptance condition \eqref{eqopti} is satisfied, there exists at least one permutation of the simulated dataset that meets the ABC criterion, namely the optimal one minimizing the distance. This permutation, denoted by \( \sigma^* = \argmin_{\sigma \in \mathcal{S}_K} d(\obsdata, \simdata_{\sigma}) \), provides a natural way to reorder the parameter vector.

In practice, the \emph{permABC} algorithm proceeds similarly to standard ABC: we sample \( \parprop \sim \pi(\cdot) \) and simulate data \( \simdata \sim f(\cdot \mid \parprop) \). If the minimum distance over all permutations satisfies the tolerance condition, we accept and return the pair \( (\parprop_{\sigma^*}, \simdata_{\sigma^*}) \). This induces a new pseudo-posterior distribution, denoted \( \pi_{\varepsilon}^*(\cdot, \cdot \mid \obsdata) \), which corresponds to the pushforward of the distribution \( \tilde{\pi}_\varepsilon \) defined in \eqref{eqpermabc_tilde} under the optimal alignment map \( \proj \):

\begin{equation}
\label{eqproj}
\pi^{*}_{\varepsilon}(\cdot, \cdot \mid \obsdata)
:= \projsharp \tilde{\pi}_{\varepsilon}(\cdot, \cdot \mid \obsdata)
\end{equation}
where
\[
\proj: (\param, \simdata) \mapsto (\param_{\sigma^*}, \simdata_{\sigma^*}).
\]
By construction, \( \proj \circ \proj = \proj \), so \( \proj \) acts as a projection. Moreover, its preimage satisfies \( \proj^{-1}(\{\param, \simdata\}) = \{(\param_{\sigma}, \simdata_{\sigma}) \mid \sigma \in \mathcal{S}_K\} \).

As a result, the projected pseudo-posterior satisfies:
\begin{equation*}
    \begin{split}
        \pi_{\varepsilon}^* (\param, \simdata \mid \obsdata) &  \propto \Tilde \pi_{\varepsilon} (\proj^{-1}(\{(\param, \simdata)\}) \mid \obsdata)\cdot \mathbb I_{\text{Im}(\proj)}(\param,\simdata)\\
        & \propto \sum_{\sigma} \Tilde \pi_{\varepsilon} (\param_{\sigma}, \simdata_{\sigma} \mid \obsdata)\cdot \mathbb I_{\text{Im}(\proj)}(\param,\simdata)\\
        % & \propto \sum_{\sigma} \pi(\param_{\sigma}) f(\simdata_{\sigma} \mid \param_{\sigma}) \mathbb{I}_{\tilde A_{\obsdata,\varepsilon}}(\simdata_{\sigma})\cdot \mathbb I_{\text{Im}(\proj)}(\param,\simdata)\\
        % & \propto \pi(\param) f(\simdata \mid \param) \mathbb{I}_{\tilde A_{\obsdata,\varepsilon}}(\simdata)\cdot \mathbb I_{\text{Im}(\proj)}(\param,\simdata)\\
        & \propto \tilde \pi_{\varepsilon}(\param, \simdata\mid \obsdata) \cdot \mathbb I_{\text{Im}(\proj)}(\param,\simdata).
    \end{split}
\end{equation*}


\noindent \old{which enforces a constraint that accepted samples must lie in the image of the projection, i.e., must be optimally aligned with the observed data.

This additional constraint effectively restricts the integration domain to a subset of datasets that are optimally ordered with respect to \( \obsdata \). We define this constrained space as}\new{This restricts accepted samples to the image of \( \proj \). We define the constrained space as}
\[
\dataspace_{\obsdata}^{*K} = \left\{ \simdata \in \dataspace^K \mid d(\obsdata, \simdata) = \min_{\sigma \in \mathcal{S}_K} d(\obsdata, \simdata_{\sigma}) \right\}.
\]
Consequently, the marginal pseudo-posterior can be expressed as:
\begin{equation}
\label{eqpermabc_star}
\pi_{\varepsilon}^*(\param \mid \obsdata) \propto \pi(\param) \int_{A_{\obsdata,\varepsilon}^*} f(\simdata \mid \param) \, d\simdata,
\end{equation}
where the new acceptance region is 
\[
A_{\obsdata,\varepsilon}^* = B_{\varepsilon}(\obsdata) \cap \dataspace_{\obsdata}^{*K}.
\]

The \emph{permABC} algorithm (see Algorithm~\ref{alg:permABC_star}) therefore defines a different approximation of the posterior distribution than classical ABC methods. In particular, it yields a different posterior approximation  \( \pi_{\varepsilon}^*(\cdot \mid \obsdata) \), which we shall compare to the vanilla ABC approximation \( \pi_{\varepsilon}(\cdot \mid \obsdata) \).


\begin{algorithm}
{\small
\SetKwInOut{Inputt}{Input}
\SetKwInOut{Outputt}{Output}
\caption{permABC Vanilla}\label{alg:permABC_star}
\textbf{Input: }observed dataset $\obsdata$, number of iterations $N$, threshold $\varepsilon>0$.\\
\textbf{Output: }a sample $(\param^{(1)}, \dots, \param^{(N)})$ from the pseudo-posterior distribution $\pi_{\varepsilon}^*(\cdot\mid \obsdata)$.\\

\BlankLine

\For{$i = 1,\dots, N$}{
    \Repeat{$d(\obsdata, \simdata_{\sigma^*}) \leq \eps$}{
        $\parprop \sim \pi(\cdot)$\;
        $\simdata \sim f(\cdot \mid \parprop)$\;
        $\sigma^* = \argmin_{\sigma \in \mathcal{S}_K} d(\obsdata, \simdata_{\sigma})$\;
    }
    $\param^{(i)} = \param'_{\sigma^*}\;$
}
}% end \small
\end{algorithm}

\subsection{Comparison with ABC and asymptotic perspective}
\label{sec:comparison}

The permutation step in permABC has two distinct asymptotic interpretations, depending on whether one focuses on the global parameter \(\globparam\) or on the labelled local parameters \(\locparam_{1:K}\).

For the global parameter \(\globparam\), once the labels of the local components are ignored, the acceptance rule of permABC falls within the same Wasserstein-type paradigm as WABC \citep{bernton2019approximate}. In our hierarchical setting, the number of compartments \(K\) plays the role of the sample size in WABC, since the comparison is performed through the empirical distributions of the observed and simulated compartments. This viewpoint is useful because it allows us to directly invoke the asymptotic intuition developed for WABC: when \(\varepsilon \to 0\) with the observations kept fixed, the pseudo-posterior approaches the exact posterior, while as the sample size increases one obtains posterior concentration under suitable identifiability conditions. In our context, this supports the view that the permutation-invariant marginal pseudo-posterior in \(\globparam\) inherits the same qualitative behaviour as \(K\) increases. Moreover, when \(\varepsilon\) is kept fixed, the resulting target should be interpreted as a coarsened posterior rather than as a Dirac mass, in the spirit of \citet{miller2019robust}. 

For the labelled local parameters \(\locparam_{1:K}\), however, the relevant comparison is no longer with WABC but with standard ABC. Indeed, permABC differs from standard ABC in one key aspect: instead of retaining all permutations that satisfy the ABC criterion, it keeps only the optimal one, namely the permutation that minimizes the discrepancy defined in Equation~\eqref{eqopti}. This choice is computationally attractive, but it alters the pseudo-posterior by constraining the pseudo-data to lie in \(\dataspace^{*K}_{\obsdata}\). A stratified variant recovering the standard pseudo-posterior is described in Appendix~\ref{sec:appendix_strat}. We now show that, for small enough \(\varepsilon\), this difference vanishes.

\begin{lemma}
\label{lempermABC_star}
Let \(\obsdata \in \dataspace^K\) denote the observed dataset, and define
\[
\varepsilon^* := \frac{1}{2}\min_{\sigma \neq \mathrm{Id}} d(\obsdata,\obsdata_\sigma).
\]
Then, for any \(\varepsilon < \varepsilon^*\) and any simulated dataset \(\simdata \in \mathbb{R}^{K \times n}\), at most one permutation satisfies the ABC criterion:
\[
\mathrm{Card}\left\{\sigma \in \mathcal{S}_K : d(\obsdata,\simdata_\sigma)\le \varepsilon\right\}\le 1.
\]
\end{lemma}

The proofs of the preceding lemma, the following theorem, and the subsequent corollary are deferred to Appendix~\ref{appendix_proofs}.

We assume that no two compartments have exactly the same observed data. This implies that \(\varepsilon^*>0\), and this condition is satisfied almost surely for continuous data. Under this assumption, permABC and standard ABC coincide for sufficiently small \(\varepsilon\), as stated below.

\begin{theorem}
\label{th:permABC_star}
Let \(\pi^*_{\varepsilon}(\cdot\mid\obsdata)\) denote the pseudo-posterior distribution resulting from permABC, and \(\pi_{\varepsilon}(\cdot\mid\obsdata)\) that obtained from standard ABC. Then, for all \(\varepsilon<\varepsilon^*\),
\[
\pi^*_{\varepsilon}(\cdot\mid\obsdata)=\pi_{\varepsilon}(\cdot\mid\obsdata).
\]
\end{theorem}

This equivalence ensures that permABC inherits the same small-tolerance asymptotic behaviour as standard ABC for the full labelled parameter vector. In particular, any consistency result for ABC immediately transfers to permABC in the regime \(\varepsilon\to0\).

\begin{corollary}
\label{col:permabc}
If the pseudo-posterior distribution \(\pi_{\varepsilon}(\cdot\mid\obsdata)\) obtained by standard ABC converges to the true posterior distribution as \(\varepsilon\to0\), then the pseudo-posterior distribution \(\pi^*_{\varepsilon}(\cdot\mid\obsdata)\) obtained by permABC also converges to the true posterior distribution as \(\varepsilon\to0\).
\end{corollary}

Thus, the efficiency gains provided by permABC come at no cost to asymptotic consistency. More generally, beyond mere consistency, one may appeal to the asymptotic theory developed for ABC by \citet{frazier2018asymptotic}, which gives general results on posterior concentration, limiting posterior shape, and asymptotic distributions under suitable regularity, identifiability, and joint \((n,\varepsilon)\)-asymptotics. In the present hierarchical setting, these results apply to permABC in the regime \(\varepsilon<\varepsilon^*\), where permABC and ABC coincide, while the WABC viewpoint described above provides the appropriate asymptotic interpretation for the marginal pseudo-posterior in \(\globparam\).